% GENERATED FILE. Edit canonical lessons, then run npm run generate:books.
% canonical-source-hash: fnv1a64:12df739278b25302
% canonical-lessons: PA-C194-identity-return, PA-C194-vidiarthi, PA-C194-vidiarthi-review, PA-C194-student-frame, PA-C194-payoff

\chapter{He or She Is a Student}
\label{ch:pa-a2-vidiarthi}
\hlchaptermodality{194}

\begin{chapteropening}
\textbf{By the end of this chapter:} \emph{I can recognise and use the familiar noun for a student in one bounded identity statement.}

Recognise and produce the student noun in one bounded identity statement through listening, speaking, reading, and writing.
\end{chapteropening}

\section[\texorpdfstring{\pa{ਉਹ} · \pa{ਹੈ} (oh; hai)}{ਉਹ · ਹੈ (oh; hai)}]{Return to a familiar identity frame}

\label{lesson:PA-C194-identity-return}



\begin{quote}
\emph{Hear:} \emph{oh; hai}; YOU SAY: \textbf{\pa{ਉਹ}; \pa{ਹੈ}}
\end{quote}

\begin{grammarlens}[title={Grammar Lens}]
Keep the familiar pronoun first and the familiar copula last. The middle slot will stay empty until the next lesson owns its new noun.
\end{grammarlens}

\subsection*{Your turn}
\begin{itemize}
\raggedright
  \item \emph{Hear:} \emph{oh}; YOU SAY: \textbf{\pa{ਉਹ}}
  \item \emph{Hear:} \emph{hai}; YOU SAY: \textbf{\pa{ਹੈ}}
  \item Say the first word, pause for the empty middle, then say the last.
\end{itemize}

\begin{tcolorbox}[breakable,colback=teal!4,colframe=teal!35!black,title={Before you move on}]
Produce the two familiar words around the empty middle slot. This opener owns no new noun.
\end{tcolorbox}
\section[\texorpdfstring{\pa{ਵਿਦਿਆਰਥੀ} (vidiarthii)}{ਵਿਦਿਆਰਥੀ (vidiarthii)}]{One noun for a student}

\label{lesson:PA-C194-vidiarthi}



\begin{quote}
Retrieve the familiar first and last words of the identity frame.
\end{quote}

\begin{grammarlens}[title={Grammar Lens: direct owner}]
\begin{quote}
\textbf{\pa{ਵਿਦਿਆਰਥੀ}} — \emph{vidiārthī} — student
\end{quote}

Hear and say the whole word before reading it. Keep it whole rather than inventing smaller pieces. This lesson owns the noun alone; a complete identity statement comes only after a separate retrieval.
\end{grammarlens}

\subsection*{Your turn}
\begin{itemize}
\raggedright
  \item \emph{Hear:} \emph{vidiārthī}; YOU SAY: \textbf{\pa{ਵਿਦਿਆਰਥੀ}}
  \item \emph{Read it:} \textbf{\pa{ਵਿਦਿਆਰਥੀ}}
  \item \emph{Write it:} copy \textbf{\pa{ਵਿਦਿਆਰਥੀ}} once, checking the whole word
\end{itemize}

\begin{tcolorbox}[breakable,colback=teal!4,colframe=teal!35!black,title={Before you move on}]
Produce the noun for “student” without looking.
\end{tcolorbox}
\section[\texorpdfstring{\pa{ਵਿਦਿਆਰਥੀ} (vidiarthii)}{ਵਿਦਿਆਰਥੀ (vidiarthii)}]{Retrieve the new noun on its own}

\label{lesson:PA-C194-vidiarthi-review}



\begin{quote}
\emph{Hear:} \emph{vidiārthī}; YOU SAY: \textbf{\pa{ਵਿਦਿਆਰਥੀ}}
\end{quote}

\subsection*{Your turn}
\begin{itemize}
\raggedright
  \item \emph{Hear:} \emph{student}; YOU SAY: \textbf{\pa{ਵਿਦਿਆਰਥੀ}}
  \item \emph{Read it:} \textbf{\pa{ਵਿਦਿਆਰਥੀ}}
  \item \emph{Write it:} cover and write \textbf{\pa{ਵਿਦਿਆਰਥੀ}}
\end{itemize}

\begin{grammarlens}[title={Grammar Lens}]
Keep one form and one meaning here. Number changes and case changes stay outside this bounded lesson.
\end{grammarlens}

\begin{tcolorbox}[breakable,colback=teal!4,colframe=teal!35!black,title={Before you move on}]
Say and then write the noun once from memory.
\end{tcolorbox}
\section[\texorpdfstring{\pa{ਉਹ} \pa{ਵਿਦਿਆਰਥੀ} \pa{ਹੈ} (oh vidiarthii hai)}{ਉਹ ਵਿਦਿਆਰਥੀ ਹੈ (oh vidiarthii hai)}]{Put the noun into one identity statement}

\label{lesson:PA-C194-student-frame}



\begin{quote}
Retrieve \textbf{\pa{ਉਹ}}, then \textbf{\pa{ਵਿਦਿਆਰਥੀ}}, then \textbf{\pa{ਹੈ}}.
\end{quote}

\begin{grammarlens}[title={Grammar Lens}]
\begin{quote}
\textbf{\pa{ਉਹ} \pa{ਵਿਦਿਆਰਥੀ} \pa{ਹੈ}} — \emph{oh vidiārthī hai} — he or she is a student
\end{quote}

Punjabi \textbf{\pa{ਉਹ}} is gender-neutral in this frame; the pronoun itself does not mark “he” versus “she.” Keep the noun in the middle and the copula last.
\end{grammarlens}

\subsection*{Your turn}
\begin{itemize}
\raggedright
  \item \emph{Hear:} \emph{oh vidiārthī hai}; YOU SAY: \textbf{\pa{ਉਹ} \pa{ਵਿਦਿਆਰਥੀ} \pa{ਹੈ}}
  \item \emph{Read it:} \textbf{\pa{ਉਹ} \pa{ਵਿਦਿਆਰਥੀ} \pa{ਹੈ}}
  \item \emph{Write it:} cover and write \textbf{\pa{ਉਹ} \pa{ਵਿਦਿਆਰਥੀ} \pa{ਹੈ}}
\end{itemize}

\begin{tcolorbox}[breakable,colback=teal!4,colframe=teal!35!black,title={Before you move on}]
Produce the complete identity statement.
\end{tcolorbox}
\section[\texorpdfstring{\pa{ਉਹ} \pa{ਵਿਦਿਆਰਥੀ} \pa{ਹੈ} (oh vidiarthii hai)}{ਉਹ ਵਿਦਿਆਰਥੀ ਹੈ (oh vidiarthii hai)}]{A bounded student-noun payoff}

\label{lesson:PA-C194-payoff}



\begin{quote}
Retrieve the noun alone before rebuilding its sentence.
\end{quote}

\subsection*{Your turn: four skills}
\begin{itemize}
\raggedright
  \item \emph{Hear:} \emph{oh vidiārthī hai}; YOU SAY: \textbf{\pa{ਉਹ} \pa{ਵਿਦਿਆਰਥੀ} \pa{ਹੈ}}
  \item \emph{Read it:} \textbf{\pa{ਉਹ} \pa{ਵਿਦਿਆਰਥੀ} \pa{ਹੈ}}
  \item \emph{Write it:} cover the line and write \textbf{\pa{ਉਹ} \pa{ਵਿਦਿਆਰਥੀ} \pa{ਹੈ}}
\end{itemize}

\begin{tcolorbox}[breakable,colback=teal!4,colframe=teal!35!black,title={Before you move on}]
Say the noun, then produce the complete identity statement.

This is the seventeenth prerequisite tranche exposed by Punjabi A2 mock drafting. It does not unblock Punjabi A2 mock 1 by itself. It is not an A2 mock, not A2 readiness or attainment, and not external or human validation.
\end{tcolorbox}
